\newpage
\section{Sformułowanie problemu i wymagania}

Cel pracy możemy sformułować jako rozwiązanie następujących problemów:
\begin{enumerate}
    \item \textbf{Problem transmisji danych na duże odległości bez infrastruktury zewnętrznej.}
          Zapewnienie niezawodnej wymiany krótkich pakietów między ruchomymi węzłami
          w obszarach bez zasięgu sieci komórkowej, przy ograniczonym budżecie.
    \item \textbf{Problem optymalizacji zajętości pasma,} co sprowadza się do minimalizacji liczby
          wysyłanych pakietów. Oznacza to opracowanie metod oszczędzania wysyłki lokalizacji GPS,
          np.\ w sytuacjach kiedy węzły są w bezruchu albo daleko od granicy strefy. Warto zauważyć,
          że rozwiązanie tego problemu jest również przydatne ze względu na oszczędzanie baterii
          niewielkich węzłów radiowych.
    \item \textbf{Problem integracji dwóch systemów.} Ponieważ interfejs użytkownika znajduje się
          w smartfonie, konieczne będzie stworzenie integracji aplikacji mobilnej z siecią radiową.
\end{enumerate}

\subsection*{Przeznaczenie systemu i charakterystyka użytkownika}

System ma służyć do monitorowania położenia zwierzęcia w czasie rzeczywistym w obszarach
pozbawionych pokrycia sieci komórkowej. Użytkownikiem docelowym jest właściciel zwierzęcia,
przebywający na działce lub w terenie otwartym - osoba niekoniecznie posiadająca wiedzę techniczną
z zakresu sieci radiowych czy programowania. Z charakterystyki użytkownika wynikają następujące
założenia projektowe:
\begin{itemize}
    \item interfejs aplikacji mobilnej musi być przejrzysty i niewymagający konfiguracji
          przy każdym uruchomieniu,
    \item konfiguracja sprzętu radiowego powinna być maksymalnie jednorazowa oraz nie powinna
          wymagać interwencji użytkownika w trakcie normalnej eksploatacji,
    \item powiadomienia o przekroczeniu strefy muszą być jednoznaczne i natychmiastowe, również
          gdy aplikacja jest zminimalizowana albo telefon ma wygaszony ekran.
\end{itemize}

Typowy scenariusz użycia wygląda następująco: właściciel przebywa na działce, zwierzę porusza się
swobodnie po terenie. Właściciel uruchamia aplikację, obserwuje pozycję zwierzęcia na mapie i otrzymuje
alarm wibracyjny lub dźwiękowy w momencie, gdy zwierzę oddali się poza zdefiniowaną strefę
bezpieczeństwa.


Tak zdefiniowany problem i charakterystyka użytkownika prowadzą do następujących wymagań funkcjonalnych:
\begin{itemize}
    \item wyświetlanie pozycji wszystkich aktywnych węzłów na mapie w czasie rzeczywistym
          z historią z bieżącej sesji,
    \item definiowanie jednej lub więcej stref bezpieczeństwa w postaci wielokątów,
          przechowywanych między sesjami,
    \item wykrywanie przekroczenia granicy strefy i powiadamianie użytkownika,
    \item adaptacyjna częstotliwość wysyłania pozycji zależna od stanu ruchu węzła,
    \item przypisanie nazwy i ikony do każdego węzła.
\end{itemize}

Dodatkowo zebrano wymagania parametryczne opisujące mierzalne cechy części składowych systemu.

\paragraph{Parametry sieci radiowej}
Węzły powinny komunikować się na minimalną odległość 2 kilometrów w terenie polnym i częściowo zalesionym.
Wymaganie to wynika z analizy scenariusza użytkowania: pies myśliwski może biec z prędkością
około 30~km/h~$\approx$~8,3~m/s. Przyjmując typowy czas od zauważenia ucieczki psa
do podjęcia reakcji przez właściciela $t = 3~\text{min} = 180~\text{s}$, pies może
oddalić się na odległość:
\begin{equation}
     d = v \cdot t = 8{,}3~\text{m/s} \cdot 180~\text{s} \approx 1500~\text{m}
\end{equation}
Wartość 2~km stanowi zatem dolne ograniczenie wymaganego zasięgu komunikacji,
zapewniające ciągłość monitorowania.
Dla porównania, komercyjne systemy śledzenia psów myśliwskich opisane w sekcji
przeglądu istniejących rozwiązań osiągają zasięgi od 14~km (Dogtra Pathfinder~2)
do 20~km (Dogtrace DOG GPS X25), co potwierdza zasadność przyjętego wymagania
jako wartości minimalnej, realistycznej do osiągnięcia w sieci mesh opartej na LoRa.

 Opóźnienie dostarczania pakietów w sieci nie powinno przekraczać 3 sekund.
Wymaganie to wynika z analizy dopuszczalnego błędu pozycji w momencie wyzwolenia alarmu
geofencingu. Przy prędkości poruszania się psa $v \approx 8{,}3$~m/s, opóźnienie
$\Delta t = 3$~s oznacza, że w chwili odebrania alarmu pies może znajdować się już
w odległości:
\begin{equation}
d = v \cdot \Delta t = 8{,}3~\text{m/s} \cdot 3~\text{s} \approx 25~\text{m}
\end{equation}
poza granicą zdefiniowanej strefy. Wartość ta jest akceptowalna, zakładając że strefy
bezpieczeństwa definiowane są z pewnym marginesem. Dla porównania, system Dogtra
Pathfinder~2 dostarcza aktualizacje pozycji co 2 sekundy, co stanowi wiarygodny punkt odniesienia


\paragraph{Parametry węzła mobilnego (obroża)}

Węzeł mobilny powinien być możliwie mały i lekki, aby zapewnić możliwość przymocowania
do obroży średniej wielkości psa (ok.\ 20~kg), w taki sposób, aby nie powodował dyskomfortu ani ograniczeń ruchowych.
Waga węzła mobilnego nie powinna przekraczać 50~g, a wymiary powinny być zbliżone do 
standardowych urządzeń umieszczanych na obrożach dla psów (np. lampki, czujniki aktywności),\ około 10~cm x 5~cm x 3~cm.

Ponadto węzeł powinien spełniać wymagania ochrony na poziomie co najmniej IP65
zgodnie z normą IEC~60529, co oznacza pełną szczelność pyłową oraz odporność
na strugi wody padające z dowolnego kierunku. Poziom ten jest wystarczający
dla warunków terenowych, w których urządzenie narażone jest na deszcz,
wilgotną roślinność oraz błoto, lecz nie na zanurzenie w wodzie.



\section{Przegląd istniejących rozwiązań}

Środowisko myśliwskie zna wiele rozwiązań na utrzymywanie psów w ściśle zdefiniowanych obszarach.
Rozważając jedynie elektroniczne rozwiązania, przedstawiono tutaj jedynie te najlepiej dopracowane.

\subsection*{Dogtra Pathfinder 2}

Jest to jeden z najbardziej zaawansowanych produktów z tej dziedziny (rysunek~\ref{fig:dogtra}). Zestaw składa się
z odbiornika, który łączy się z telefonem przez Bluetooth, oraz jednego lub więcej (maksymalnie 21)
nadajników umieszczonych na obrożach elektronicznych. Urządzenia łączą się ze sobą po falach
radiowych. Odbiornik zarówno przekazuje dane z nadajników, jak też daje możliwość użycia przycisku,
który może być zaprogramowany do wzbudzania impulsów elektrycznych w obroży, emitowania światła,
dźwięku albo wibracji w nadajnikach. Aplikacja \textit{PATHFINDER2} łącząca się przez Bluetooth
z odbiornikiem służy do wyświetlania dokładnej pozycji wszystkich skonfigurowanych nadajników,
podglądania ich stanu, sprawdzania naładowania baterii nadajników oraz statusu połączenia z GPS.
Pozatym aplikacja umożliwia nagrywanie lokalizacji w sesje do późniejszego odtworzenia, nawigację
na Mapach Google, pomiary odległości oraz sterowanie modułami na nadajnikach (wibracja, dźwięk,
światło, impuls elektryczny).

Ciekawą funkcją jest możliwość definiowania różnego rodzaju ,,E-Ogrodzeń''. Firma Dogtra opracowała
trzy rodzaje:
\begin{enumerate}
    \item \textbf{Mobile-Fence} - ogrodzenie poruszające się w oparciu o lokalizację smartfona
          i wysyłające powiadomienia.
    \item \textbf{Geo-Fence} - konfiguracja statycznego ogrodzenia wysyłającego powiadomienia.
    \item \textbf{E-Fence} - statyczna konfiguracja ogrodzenia wysyłająca psu automatyczną
          korektę i powiadomienie.\footnote{Dogtra Pathfinder 2 - instrukcja obsługi, s.~26.
          Dostępna pod adresem:
          \url{https://i00.eu/file/403/15309-manualplpathfinder2.pdf}}
\end{enumerate}

Producent deklaruje następujące parametry pracy:
\begin{itemize}
    \item zasięg do 9 mil (ok.\ 14~484~m),
    \item aktualizacje pozycji co 2 sekundy,
    \item użytkowanie bez dodatkowych opłat abonamentowych.
\end{itemize}

% TODO: Skopiuj pliki dogtra.jpg oraz dogtraremote.png do katalogu img/
%       Plik dogtraremote.webp wymaga konwersji do formatu PNG przed użyciem.
\begin{figure}[!h]
    \centering
    \begin{subfigure}{0.48\linewidth}
        \centering
        \includegraphics[width=\linewidth]{dogtra.jpg}
        \caption{Odbiornik}
        \label{fig:dogtra-receiver}
    \end{subfigure}
    \hfill
    \begin{subfigure}{0.48\linewidth}
        \centering
        \includegraphics[width=\linewidth]{dogtraremote.png}
        \caption{Pilot zdalnego sterowania}
        \label{fig:dogtra-remote}
    \end{subfigure}
    \caption{Dogtra Pathfinder 2}
    \label{fig:dogtra}
\end{figure}

\subsection*{Dogtrace DOG GPS X25}

Jedynie nieznacznie różnym podejściem wykazała się firma Dogtrace, która w swoim rozwiązaniu
zdecydowała się zrezygnować z konieczności użycia smartfona. X25 składa się z identycznego
funkcjonalnie nadajnika oraz znacznie rozbudowanego odbiornika - wielkości mniejszego smartfona -
który na wbudowanym wyświetlaczu LCD pokazuje kierunek i odległość do nadajnika (rysunek~\ref{fig:dogtrace-x25}). Czeska firma
udostępniła lepszej jakości informacje o funkcjonowaniu swojego produktu. Z oferty na stronie
internetowej\footnote{\url{https://www.obroza-elektryczna.pl/dogtrace-dog-gps-x25}} dowiadujemy
się, że:
\begin{itemize}
    \item możliwe jest śledzenie do 19 psów jednocześnie,
    \item bateria nadajnika i odbiornika powinna wystarczyć na ponad 40 godzin pracy (Li-Pol
          1900~mAh),
    \item pozycjonowanie modułów zapewnia skorzystanie z kilku systemów: GPS, GALILEO, GLONASS,
          a nadajnik z odbiornikiem łączą się po LoRa,
    \item deklarowany zasięg to 20~km dla komunikacji ,,w linii wzroku'',
    \item pełna zanurzalność nadajnika i odbiornika.
\end{itemize}

Dodatkowo dostępne są funkcje takie jak: KOMPAS (kierunek na północ magnetyczną), BEEPER
(wykrywanie ruchu lub bezruchu psa), FENCE (okrągły płot - akustyczna granica wyznaczająca
przestrzeń dla psa), WAYPOINT (możliwość zapisania 4 współrzędnych GPS odbiornika i nawigacja
do tych punktów), tryb CAR (umożliwiający korzystanie z odbiornika w pojeździe) oraz sygnał
dźwiękowy zastępujący gwizdek (3 różne dźwięki do wyboru).

% TODO: Skopiuj plik Dogtrace.png do katalogu img/
\begin{figure}[!h]
    \label{fig:dogtrace-x25}
    \centering
    \includegraphics[width=0.5\linewidth]{Dogtrace.png}
    \caption{Dogtrace DOG GPS X25}
\end{figure}

\subsection*{Porównanie z proponowanym rozwiązaniem}

Oba opisane systemy reprezentują podejście typowe dla komercyjnych rozwiązań myśliwskich:
dedykowany sprzęt, bezpośrednia komunikacja nadajnik-odbiornik oraz wysoka cena zakupu.
Proponowany w niniejszej pracy system różni się od nich pod trzema istotnymi względami.

Po pierwsze, komunikacja opiera się na sieci typu mesh, w której każdy węzeł może pełnić
rolę przekaźnika. Oznacza to, że zasięg systemu nie jest ograniczony odległością między
nadajnikiem a odbiornikiem, lecz może być rozszerzany przez dodawanie kolejnych węzłów
pośredniczących - co nie jest możliwe w architekturze punkt-punkt stosowanej przez Dogtrę
i Dogtrace.

Po drugie, system zbudowany jest z ogólnodostępnych komponentów elektronicznych, co 
przekłada się na znacząco niższy koszt w porównaniu z dedykowanym sprzętem komercyjnym.

Po trzecie, liczba węzłów sieci jest nieograniczona przez architekturę systemu, podczas gdy
rozwiązania komercyjne narzucają sztywny limit (21 węzłów dla Dogtra, 19 dla Dogtrace).

Zestawienie kluczowych cech wszystkich trzech rozwiązań przedstawia tabela~\ref{tab:porownanie}.

\begin{table}[!h]
    \caption{Porównanie systemów śledzenia psów myśliwskich}
    \label{tab:porownanie}
    \centering
    \begin{tabular}{|l|c|c|c|}
        \hline
        \textbf{Cecha} & \textbf{Dogtra Pathfinder 2} & \textbf{Dogtrace X25} & \textbf{Proponowany system} \\
        \hline
        Architektura sieci      & punkt-punkt & punkt-punkt & mesh \\
        Zasięg                  & $\sim$14~km & $\sim$20~km & $\geq$2~km\footnotemark[1] \\
        Maks. liczba węzłów     & 21          & 19          & nieograniczona\footnotemark[2] \\
        Aktualizacja pozycji    & co 2~s      & co 3~s      & $\leq$3~s \\
        Geofencing              & tak         & tak         & tak \\
        Kształt strefy          & dowolny wielokąt & okrąg  & dowolny wielokąt \\
        Ochrona                 & IPX9K       & IP68        & IP65 \\
        Koszt                   & 2500 PLN    & 2300 PLN    & 400 PLN\footnotemark[3] \\
        \hline
    \end{tabular}
\end{table}
\footnotetext[1]{Topologia mesh pozwala na rozszerzanie zasięgu przez rozmieszczenie dodatkowych węzłów pośredniczących - podana wartość dotyczy komunikacji między sąsiednimi węzłami sieci.}
\footnotetext[2]{Nieograniczona z punktu widzenia architektury systemu - w praktyce ogranicza ją przepustowość współdzielonego kanału radiowego: wraz ze wzrostem liczby węzłów rośnie ruch w sieci, co może prowadzić do wzrostu zajętości kanału (ograniczonej prawnie).}
\footnotetext[3]{+ 200~PLN za każdy kolejny węzeł.}
